% normalization.tex — BN / LN / IN / GN / RMSNorm side-by-side.
% Используется в lectures/03-mlp-norm и lectures/04-cnn (как припоминание).
% Идея визуализации: тензор [N, C, H*W] изображён как сетка кубиков; цветом
% выделена область, по которой считаются среднее и дисперсия.
%
% Скомпилировать: python figures/render.py figures/tikz/normalization.tex
% Подключить в .qmd: ![](../figures/tikz/_build/normalization.svg)

\documentclass[tikz,border=4pt]{standalone}
% Общая преамбула для всех TikZ-фигур курса ml3-neo.
% Подключается:
%   - render.py при компиляции через Tectonic в standalone SVG/PDF;
%   - Quarto _quarto.yml через header-includes для PDF-сборки лекций;
%   - лекциями вручную через % Общая преамбула для всех TikZ-фигур курса ml3-neo.
% Подключается:
%   - render.py при компиляции через Tectonic в standalone SVG/PDF;
%   - Quarto _quarto.yml через header-includes для PDF-сборки лекций;
%   - лекциями вручную через % Общая преамбула для всех TikZ-фигур курса ml3-neo.
% Подключается:
%   - render.py при компиляции через Tectonic в standalone SVG/PDF;
%   - Quarto _quarto.yml через header-includes для PDF-сборки лекций;
%   - лекциями вручную через \input{figures/tikz/_preamble.tex}.
%
% Все цвета определены через xcolor чтобы хорошо смотреться и в light, и в dark теме.

\usepackage{tikz}
\usepackage{xcolor}
\usepackage{amsmath, amssymb}
\usepackage{bm}

\usetikzlibrary{
  positioning,
  arrows.meta,
  calc,
  fit,
  shapes.geometric,
  shapes.misc,
  backgrounds,
  decorations.pathreplacing,
  decorations.pathmorphing,
  patterns,
  matrix,
  3d
}

% --- Цветовая палитра курса ----------------------------------------------
% Совместима с акцентом сайта (#4f46e5) и нейтральная по отношению к темам.
\definecolor{ml3accent}{HTML}{4F46E5}
\definecolor{ml3accent2}{HTML}{0EA5E9}
\definecolor{ml3warm}{HTML}{F59E0B}
\definecolor{ml3rose}{HTML}{E11D48}
\definecolor{ml3green}{HTML}{10B981}
\definecolor{ml3muted}{HTML}{6B7280}
\definecolor{ml3soft}{HTML}{E5E7EB}
\definecolor{ml3ink}{HTML}{111827}

% --- Общие стили nodes/arrows --------------------------------------------
\tikzset{
  >={Stealth[round]},
  ml3/box/.style    = {draw=ml3ink, line width=0.7pt, rounded corners=2pt,
                       fill=ml3soft, minimum size=10mm, font=\small},
  ml3/ghost/.style  = {draw=ml3muted, dashed, rounded corners=2pt,
                       inner sep=4pt, font=\footnotesize\itshape\color{ml3muted}},
  ml3/op/.style     = {draw=ml3accent, fill=white, line width=0.8pt,
                       circle, minimum size=6mm, inner sep=0pt, font=\small},
  ml3/tensor/.style = {draw=ml3accent2, line width=0.6pt, fill=ml3accent2!12,
                       rounded corners=1pt, minimum height=8mm, font=\footnotesize},
  ml3/arrow/.style  = {->, line width=0.7pt, draw=ml3ink},
  ml3/dashedarrow/.style = {->, line width=0.6pt, draw=ml3muted, dashed},
  ml3/label/.style  = {font=\footnotesize\color{ml3muted}}
}
.
%
% Все цвета определены через xcolor чтобы хорошо смотреться и в light, и в dark теме.

\usepackage{tikz}
\usepackage{xcolor}
\usepackage{amsmath, amssymb}
\usepackage{bm}

\usetikzlibrary{
  positioning,
  arrows.meta,
  calc,
  fit,
  shapes.geometric,
  shapes.misc,
  backgrounds,
  decorations.pathreplacing,
  decorations.pathmorphing,
  patterns,
  matrix,
  3d
}

% --- Цветовая палитра курса ----------------------------------------------
% Совместима с акцентом сайта (#4f46e5) и нейтральная по отношению к темам.
\definecolor{ml3accent}{HTML}{4F46E5}
\definecolor{ml3accent2}{HTML}{0EA5E9}
\definecolor{ml3warm}{HTML}{F59E0B}
\definecolor{ml3rose}{HTML}{E11D48}
\definecolor{ml3green}{HTML}{10B981}
\definecolor{ml3muted}{HTML}{6B7280}
\definecolor{ml3soft}{HTML}{E5E7EB}
\definecolor{ml3ink}{HTML}{111827}

% --- Общие стили nodes/arrows --------------------------------------------
\tikzset{
  >={Stealth[round]},
  ml3/box/.style    = {draw=ml3ink, line width=0.7pt, rounded corners=2pt,
                       fill=ml3soft, minimum size=10mm, font=\small},
  ml3/ghost/.style  = {draw=ml3muted, dashed, rounded corners=2pt,
                       inner sep=4pt, font=\footnotesize\itshape\color{ml3muted}},
  ml3/op/.style     = {draw=ml3accent, fill=white, line width=0.8pt,
                       circle, minimum size=6mm, inner sep=0pt, font=\small},
  ml3/tensor/.style = {draw=ml3accent2, line width=0.6pt, fill=ml3accent2!12,
                       rounded corners=1pt, minimum height=8mm, font=\footnotesize},
  ml3/arrow/.style  = {->, line width=0.7pt, draw=ml3ink},
  ml3/dashedarrow/.style = {->, line width=0.6pt, draw=ml3muted, dashed},
  ml3/label/.style  = {font=\footnotesize\color{ml3muted}}
}
.
%
% Все цвета определены через xcolor чтобы хорошо смотреться и в light, и в dark теме.

\usepackage{tikz}
\usepackage{xcolor}
\usepackage{amsmath, amssymb}
\usepackage{bm}

\usetikzlibrary{
  positioning,
  arrows.meta,
  calc,
  fit,
  shapes.geometric,
  shapes.misc,
  backgrounds,
  decorations.pathreplacing,
  decorations.pathmorphing,
  patterns,
  matrix,
  3d
}

% --- Цветовая палитра курса ----------------------------------------------
% Совместима с акцентом сайта (#4f46e5) и нейтральная по отношению к темам.
\definecolor{ml3accent}{HTML}{4F46E5}
\definecolor{ml3accent2}{HTML}{0EA5E9}
\definecolor{ml3warm}{HTML}{F59E0B}
\definecolor{ml3rose}{HTML}{E11D48}
\definecolor{ml3green}{HTML}{10B981}
\definecolor{ml3muted}{HTML}{6B7280}
\definecolor{ml3soft}{HTML}{E5E7EB}
\definecolor{ml3ink}{HTML}{111827}

% --- Общие стили nodes/arrows --------------------------------------------
\tikzset{
  >={Stealth[round]},
  ml3/box/.style    = {draw=ml3ink, line width=0.7pt, rounded corners=2pt,
                       fill=ml3soft, minimum size=10mm, font=\small},
  ml3/ghost/.style  = {draw=ml3muted, dashed, rounded corners=2pt,
                       inner sep=4pt, font=\footnotesize\itshape\color{ml3muted}},
  ml3/op/.style     = {draw=ml3accent, fill=white, line width=0.8pt,
                       circle, minimum size=6mm, inner sep=0pt, font=\small},
  ml3/tensor/.style = {draw=ml3accent2, line width=0.6pt, fill=ml3accent2!12,
                       rounded corners=1pt, minimum height=8mm, font=\footnotesize},
  ml3/arrow/.style  = {->, line width=0.7pt, draw=ml3ink},
  ml3/dashedarrow/.style = {->, line width=0.6pt, draw=ml3muted, dashed},
  ml3/label/.style  = {font=\footnotesize\color{ml3muted}}
}


\begin{document}
\begin{tikzpicture}[
  cube/.style={draw=ml3ink, line width=0.35pt, fill=ml3soft},
  hl/.style={draw=ml3ink, line width=0.4pt, fill=ml3accent!55},
  panel title/.style={font=\small\bfseries\color{ml3ink}, anchor=north},
  panel sub/.style={font=\scriptsize\color{ml3muted}, anchor=north},
  x=4mm, y=4mm
]

% --- Параметры тензора (визуально маленькие, иначе фигура огромная) --------
\def\Nsamples{4}    % batch
\def\Cchans{6}      % channels
\def\Spatial{4}     % H*W flattened
\def\panelW{38mm}
\def\panelGap{2mm}

% --- Макрос: одна "панель" с типом нормализации --------------------------
% #1 — anchor/coordinate
% #2 — заголовок (BatchNorm и т.п.)
% #3 — формула области (по чему усредняем)
% #4 — тип подсветки: "bn", "ln", "in", "gn", "rms"
\newcommand{\normpanel}[4]{
  \begin{scope}[shift={#1}]
    % фоновая рамка
    \node[draw=ml3muted, dashed, rounded corners=2pt, inner sep=2mm,
          minimum width=\panelW, minimum height=30mm, anchor=north west] (frame) at (0,0) {};

    % тензор: N строк × C столбцов × Spatial кубиков
    \foreach \n in {1,...,\Nsamples} {
      \foreach \c in {1,...,\Cchans} {
        \foreach \s in {1,...,\Spatial} {
          \pgfmathsetmacro\px{(\c-1)*1.0 + (\s-1)*0.18 + 1}
          \pgfmathsetmacro\py{-(\n-1)*1.0 - (\s-1)*0.18 - 4}

          % решаем подсветку
          \def\highlight{0}
          \ifnum\pdfstrcmp{#4}{bn}=0
            \ifnum\c=2 \def\highlight{1}\fi
          \fi
          \ifnum\pdfstrcmp{#4}{ln}=0
            \ifnum\n=2 \def\highlight{1}\fi
          \fi
          \ifnum\pdfstrcmp{#4}{in}=0
            \ifnum\n=2 \ifnum\c=2 \def\highlight{1}\fi\fi
          \fi
          \ifnum\pdfstrcmp{#4}{gn}=0
            \ifnum\n=2
              \ifnum\c<5
                \ifnum\c>1 \def\highlight{1}\fi
              \fi
            \fi
          \fi
          \ifnum\pdfstrcmp{#4}{rms}=0
            \ifnum\n=2 \def\highlight{1}\fi
          \fi

          \ifnum\highlight=1
            \draw[hl] (\px, \py) rectangle ++(0.85, 0.85);
          \else
            \draw[cube] (\px, \py) rectangle ++(0.85, 0.85);
          \fi
        }
      }
    }

    % оси
    \node[ml3/label, anchor=east] at (0.6, -1.5) {$N$};
    \node[ml3/label, anchor=south] at (4.0, 0.2) {$C$};
    \node[ml3/label, anchor=west] at (8.2, -0.6) {$H{\cdot}W$};

    % заголовок и подформула
    \node[panel title] at ($(frame.north) + (0, -1mm)$) {#2};
    \node[panel sub]   at ($(frame.north) + (0, -5mm)$) {#3};
  \end{scope}
}

% --- Размещение пяти панелей -----------------------------------------------
\normpanel{(0,0)}{BatchNorm}{$\mu_c, \sigma_c$ — по $N{\times}H{\times}W$}{bn}
\normpanel{(45mm,0)}{LayerNorm}{$\mu_n, \sigma_n$ — по $C{\times}H{\times}W$}{ln}
\normpanel{(90mm,0)}{InstanceNorm}{$\mu_{n,c}, \sigma_{n,c}$ — по $H{\times}W$}{in}
\normpanel{(0,-38mm)}{GroupNorm}{$\mu_{n,g}, \sigma_{n,g}$ — по группе $G$}{gn}
\normpanel{(45mm,-38mm)}{RMSNorm}{нет $\mu$; только $\mathrm{RMS}_n$}{rms}

% --- Легенда формулы -------------------------------------------------------
\node[anchor=north west, font=\footnotesize, text width=42mm,
      inner sep=2mm, draw=ml3muted, rounded corners=2pt]
  at (90mm,-38mm)
{
  \textbf{Общая форма:}\\[2pt]
  $\hat x = \dfrac{x - \mu}{\sqrt{\sigma^2 + \varepsilon}}$\\[3pt]
  $y = \gamma \hat x + \beta$\\[4pt]
  Различие — \textit{по какой оси} считаем $\mu$ и $\sigma^2$. Подсвечена область усреднения для одного элемента батча.
};

\end{tikzpicture}
\end{document}
